
\documentclass{article}
\usepackage{colortbl}
\usepackage{makecell}
\usepackage{multirow}
\usepackage{supertabular}

\begin{document}

\newcounter{utterance}

\twocolumn

{ \footnotesize  \setcounter{utterance}{1}
\setlength{\tabcolsep}{0pt}
\begin{supertabular}{c@{$\;$}|p{.15\linewidth}@{}p{.15\linewidth}p{.15\linewidth}p{.15\linewidth}p{.15\linewidth}p{.15\linewidth}}

    \# & $\;$A & \multicolumn{4}{c}{Game Master} & $\;\:$B\\
    \hline 

    \theutterance \stepcounter{utterance}  

    & & \multicolumn{4}{p{0.6\linewidth}}{\cellcolor[rgb]{0.9,0.9,0.9}{%
	\makecell[{{p{\linewidth}}}]{% 
	  \tt {\tiny [A$\langle$GM]}  
	 Du nimmst an einem kollaborativen Verhandlungspiel Teil.\\ \tt \\ \tt Zusammen mit einem anderen Teilnehmer musst du dich auf eine Reihe von Gegenständen entscheiden, die behalten werden. Jeder von euch hat eine persönliche Verteilung über die Wichtigkeit der einzelnen Gegenstände. Jeder von euch hat eine eigene Meinung darüber, wie wichtig jeder einzelne Gegenstand ist (Gegenstandswichtigkeit). Du kennst die Wichtigkeitsverteilung des anderen Spielers nicht. Zusätzlich siehst du, wie viel Aufwand jeder Gegenstand verursacht.  \\ \tt Ihr dürft euch nur auf eine Reihe von Gegenständen einigen, wenn der Gesamtaufwand der ausgewählten Gegenstände den Maximalaufwand nicht überschreitet:\\ \tt \\ \tt Maximalaufwand = 2984\\ \tt \\ \tt Hier sind die einzelnen Aufwände der Gegenstände:\\ \tt \\ \tt Aufwand der Gegenstände = {"C76": 350, "C38": 276, "C56": 616, "A03": 737, "A07": 531, "A83": 389, "C98": 24, "C08": 125, "C62": 338, "C00": 356, "C10": 143, "C81": 117, "B38": 257, "C03": 921, "C32": 789}\\ \tt \\ \tt Hier ist deine persönliche Verteilung der Wichtigkeit der einzelnen Gegenstände:\\ \tt \\ \tt Werte der Gegenstandswichtigkeit = {"C76": 138, "C38": 583, "C56": 868, "A03": 822, "A07": 783, "A83": 65, "C98": 262, "C08": 121, "C62": 508, "C00": 780, "C10": 461, "C81": 484, "B38": 668, "C03": 389, "C32": 808}\\ \tt \\ \tt Ziel:\\ \tt \\ \tt Dein Ziel ist es, eine Reihe von Gegenständen auszuhandeln, die dir möglichst viel bringt (d. h. Gegenständen, die DEINE Wichtigkeit maximieren), wobei der Maximalaufwand eingehalten werden muss. Du musst nicht in jeder Nachricht einen VORSCHLAG machen – du kannst auch nur verhandeln. Alle Taktiken sind erlaubt!\\ \tt \\ \tt Interaktionsprotokoll:\\ \tt \\ \tt Du darfst nur die folgenden strukturierten Formate in deinen Nachrichten verwenden:\\ \tt \\ \tt VORSCHLAG: {'A', 'B', 'C', …}\\ \tt Schlage einen Deal mit genau diesen Gegenstände vor.\\ \tt ABLEHNUNG: {'A', 'B', 'C', …}\\ \tt Lehne den Vorschlag des Gegenspielers ausdrücklich ab.\\ \tt ARGUMENT: {'...'}\\ \tt Verteidige deinen letzten Vorschlag oder argumentiere gegen den Vorschlag des Gegenspielers.\\ \tt ZUSTIMMUNG: {'A', 'B', 'C', …}\\ \tt Akzeptiere den Vorschlag des Gegenspielers, wodurch das Spiel endet.\\ \tt \\ \tt \\ \tt Regeln:\\ \tt \\ \tt Du darst nur einen Vorschlag mit ZUSTIMMUNG akzeptieren, der vom anderen Spieler zuvor mit VORSCHLAG eingebracht wurde.\\ \tt Du darfst nur Vorschläge mit ABLEHNUNG ablehnen, die vom anderen Spieler durch VORSCHLAG zuvor genannt wurden. \\ \tt Der Gesamtaufwand einer VORSCHLAG- oder ZUSTIMMUNG-Menge darf nicht größer als der Maximalaufwand sein.  \\ \tt Offenbare deine versteckte Wichtigkeitsverteilung nicht.\\ \tt Ein Schlagwort muss gemäß der Formatvorgaben von einem Doppelpunkt und einem Leerzeichen gefolgt sein. Das Argument muss eine Python-Menge mit 0 oder mehr Strings sein.  \\ \tt Die Form muss also SCHLAGWORT: {...} sein.\\ \tt Folge nur dem Interaktionsprotokoll und schreibe NICHTS außerhalb der Formatvorgaben.\\ \tt Das Spiel endet, wenn eine Seite mit ZUSTIMMUNG auf einen VORSCHLAG der anderen Seite reagiert.  \\ \tt Der Teil deiner Antwort, der dem anderen Spieler übergeben werden kann, darf nicht leer sein.  \\ \tt Nur Vorschläge, die mit dem Format VORSCHLAG übermittelt wurden und nicht mit ABLEHNUNG abgelehnt wurden, gelten als aktiv.  \\ \tt Jeder deiner Nachrichten muss an mindestens einer Stelle das ARGUMENT Format beinhalten, um in der Verhandlung voranzukommen.\\ \tt  
	  } 
	   } 
	   } 
	 & \\ 
 

    \theutterance \stepcounter{utterance}  

    & \multicolumn{4}{p{0.6\linewidth}}{\cellcolor[rgb]{0.8,1,0.9}{%
	\makecell[{{p{\linewidth}}}]{% 
	  \tt {\tiny [A$\rangle$GM]}  
	 VORSCHLAG: {'C98', 'C81', 'C10', 'B38', 'C00', 'C38', 'C62', 'A07', 'C56', 'C08'}\\ \tt ARGUMENT: {"Die ausgewählten Gegenstände sind optimal, um den maximalen Wert innerhalb der Aufwandsgrenze zu erzielen."} 
	  } 
	   } 
	   } 
	 & & \\ 
 

    \theutterance \stepcounter{utterance}  

    & & \multicolumn{4}{p{0.6\linewidth}}{\cellcolor[rgb]{0.9,0.9,0.9}{%
	\makecell[{{p{\linewidth}}}]{% 
	  \tt {\tiny [GM$\rangle$B]}  
	 Du nimmst an einem kollaborativen Verhandlungspiel Teil.\\ \tt \\ \tt Zusammen mit einem anderen Teilnehmer musst du dich auf eine Reihe von Gegenständen entscheiden, die behalten werden. Jeder von euch hat eine persönliche Verteilung über die Wichtigkeit der einzelnen Gegenstände. Jeder von euch hat eine eigene Meinung darüber, wie wichtig jeder einzelne Gegenstand ist (Gegenstandswichtigkeit). Du kennst die Wichtigkeitsverteilung des anderen Spielers nicht. Zusätzlich siehst du, wie viel Aufwand jeder Gegenstand verursacht.  \\ \tt Ihr dürft euch nur auf eine Reihe von Gegenständen einigen, wenn der Gesamtaufwand der ausgewählten Gegenstände den Maximalaufwand nicht überschreitet:\\ \tt \\ \tt Maximalaufwand = 2984\\ \tt \\ \tt Hier sind die einzelnen Aufwände der Gegenstände:\\ \tt \\ \tt Aufwand der Gegenstände = {"C76": 350, "C38": 276, "C56": 616, "A03": 737, "A07": 531, "A83": 389, "C98": 24, "C08": 125, "C62": 338, "C00": 356, "C10": 143, "C81": 117, "B38": 257, "C03": 921, "C32": 789}\\ \tt \\ \tt Hier ist deine persönliche Verteilung der Wichtigkeit der einzelnen Gegenstände:\\ \tt \\ \tt Werte der Gegenstandswichtigkeit = {"C56": 65, "A03": 121, "C32": 138, "A07": 262, "C00": 389, "B38": 461, "C38": 484, "C62": 508, "C81": 583, "C10": 668, "C03": 780, "C98": 783, "C76": 808, "C08": 822, "A83": 868}\\ \tt \\ \tt Ziel:\\ \tt \\ \tt Dein Ziel ist es, eine Reihe von Gegenständen auszuhandeln, die dir möglichst viel bringt (d. h. Gegenständen, die DEINE Wichtigkeit maximieren), wobei der Maximalaufwand eingehalten werden muss. Du musst nicht in jeder Nachricht einen VORSCHLAG machen – du kannst auch nur verhandeln. Alle Taktiken sind erlaubt!\\ \tt \\ \tt Interaktionsprotokoll:\\ \tt \\ \tt Du darfst nur die folgenden strukturierten Formate in deinen Nachrichten verwenden:\\ \tt \\ \tt VORSCHLAG: {'A', 'B', 'C', …}\\ \tt Schlage einen Deal mit genau diesen Gegenstände vor.\\ \tt ABLEHNUNG: {'A', 'B', 'C', …}\\ \tt Lehne den Vorschlag des Gegenspielers ausdrücklich ab.\\ \tt ARGUMENT: {'...'}\\ \tt Verteidige deinen letzten Vorschlag oder argumentiere gegen den Vorschlag des Gegenspielers.\\ \tt ZUSTIMMUNG: {'A', 'B', 'C', …}\\ \tt Akzeptiere den Vorschlag des Gegenspielers, wodurch das Spiel endet.\\ \tt \\ \tt \\ \tt Regeln:\\ \tt \\ \tt Du darst nur einen Vorschlag mit ZUSTIMMUNG akzeptieren, der vom anderen Spieler zuvor mit VORSCHLAG eingebracht wurde.\\ \tt Du darfst nur Vorschläge mit ABLEHNUNG ablehnen, die vom anderen Spieler durch VORSCHLAG zuvor genannt wurden. \\ \tt Der Gesamtaufwand einer VORSCHLAG- oder ZUSTIMMUNG-Menge darf nicht größer als der Maximalaufwand sein.  \\ \tt Offenbare deine versteckte Wichtigkeitsverteilung nicht.\\ \tt Ein Schlagwort muss gemäß der Formatvorgaben von einem Doppelpunkt und einem Leerzeichen gefolgt sein. Das Argument muss eine Python-Menge mit 0 oder mehr Strings sein.  \\ \tt Die Form muss also SCHLAGWORT: {...} sein.\\ \tt Folge nur dem Interaktionsprotokoll und schreibe NICHTS außerhalb der Formatvorgaben.\\ \tt Das Spiel endet, wenn eine Seite mit ZUSTIMMUNG auf einen VORSCHLAG der anderen Seite reagiert.  \\ \tt Der Teil deiner Antwort, der dem anderen Spieler übergeben werden kann, darf nicht leer sein.  \\ \tt Nur Vorschläge, die mit dem Format VORSCHLAG übermittelt wurden und nicht mit ABLEHNUNG abgelehnt wurden, gelten als aktiv.  \\ \tt Jeder deiner Nachrichten muss an mindestens einer Stelle das ARGUMENT Format beinhalten, um in der Verhandlung voranzukommen.\\ \tt \\ \tt \\ \tt Du erhältst nun die erste Nachricht des anderen Spielers.\\ \tt \\ \tt VORSCHLAG: {'C98', 'C81', 'C10', 'B38', 'C00', 'C38', 'C62', 'A07', 'C56', 'C08'}\\ \tt ARGUMENT: {"Die ausgewählten Gegenstände sind optimal, um den maximalen Wert innerhalb der Aufwandsgrenze zu erzielen."} 
	  } 
	   } 
	   } 
	 & \\ 
 

    \theutterance \stepcounter{utterance}  

    & & & \multicolumn{4}{p{0.6\linewidth}}{\cellcolor[rgb]{1,0.85,0.72}{%
	\makecell[{{p{\linewidth}}}]{% 
	  \tt {\tiny [GM$\langle$B]}  
	 ABLEHNUNG: {'C98', 'C81', 'C10', 'B38', 'C00', 'C38', 'C62', 'A07', 'C56', 'C08'}\\ \tt ARGUMENT: {"Die ausgewählten Gegenstände enthalten Elemente mit niedriger Priorität für mich, während hochpriorisierte Optionen wie 'A83' und 'C76' nicht berücksichtigt wurden."} 
	  } 
	   } 
	   } 
	  \\ 
 

    \theutterance \stepcounter{utterance}  

    & & \multicolumn{4}{p{0.6\linewidth}}{\cellcolor[rgb]{0.9,0.9,0.9}{%
	\makecell[{{p{\linewidth}}}]{% 
	  \tt {\tiny [A$\langle$GM]}  
	 ABLEHNUNG: {'C98', 'C81', 'C10', 'B38', 'C00', 'C38', 'C62', 'A07', 'C56', 'C08'}\\ \tt ARGUMENT: {"Die ausgewählten Gegenstände enthalten Elemente mit niedriger Priorität für mich, während hochpriorisierte Optionen wie 'A83' und 'C76' nicht berücksichtigt wurden."} 
	  } 
	   } 
	   } 
	 & \\ 
 

    \theutterance \stepcounter{utterance}  

    & \multicolumn{4}{p{0.6\linewidth}}{\cellcolor[rgb]{0.8,1,0.9}{%
	\makecell[{{p{\linewidth}}}]{% 
	  \tt {\tiny [A$\rangle$GM]}  
	 VORSCHLAG: {'C98', 'C81', 'C10', 'B38', 'C00', 'C38', 'C62', 'A07', 'C08', 'A83', 'C76'}\\ \tt ARGUMENT: {"Ich habe sowohl 'A83' als auch 'C76' aufgenommen, wie von Ihnen gewünscht. Um den Gesamtaufwand einzuhalten, musste ich 'C56' entfernen."} 
	  } 
	   } 
	   } 
	 & & \\ 
 

    \theutterance \stepcounter{utterance}  

    & & \multicolumn{4}{p{0.6\linewidth}}{\cellcolor[rgb]{0.9,0.9,0.9}{%
	\makecell[{{p{\linewidth}}}]{% 
	  \tt {\tiny [GM$\rangle$B]}  
	 VORSCHLAG: {'C98', 'C81', 'C10', 'B38', 'C00', 'C38', 'C62', 'A07', 'C08', 'A83', 'C76'}\\ \tt ARGUMENT: {"Ich habe sowohl 'A83' als auch 'C76' aufgenommen, wie von Ihnen gewünscht. Um den Gesamtaufwand einzuhalten, musste ich 'C56' entfernen."} 
	  } 
	   } 
	   } 
	 & \\ 
 

    \theutterance \stepcounter{utterance}  

    & & & \multicolumn{4}{p{0.6\linewidth}}{\cellcolor[rgb]{1,0.85,0.72}{%
	\makecell[{{p{\linewidth}}}]{% 
	  \tt {\tiny [GM$\langle$B]}  
	 assistant\\ \tt ABLEHNUNG: {'C98', 'C81', 'C10', 'B38', 'C00', 'C38', 'C62', 'A07', 'C08', 'A83', 'C76'}\\ \tt ARGUMENT: {"Die momentane Zusammensetzung könnte optimiert werden: Durch Austausch bestimmter Items mit niedrigerer Priorität gegen solche mit höherer Priorität, die insgesamt den Maximallaufwand nicht überschreiten."} 
	  } 
	   } 
	   } 
	  \\ 
 

    \theutterance \stepcounter{utterance}  

    & & \multicolumn{4}{p{0.6\linewidth}}{\cellcolor[rgb]{0.9,0.9,0.9}{%
	\makecell[{{p{\linewidth}}}]{% 
	  \tt {\tiny [GM$\rangle$B]}  
	 Deine Antwort hat eine unmarkierte Sequenz enthalten oder Du hast STRATEGIC REASONING mehr als ein mal verwendet. Du darfst nur die strukturierten Formate verwenden, die in der ersten Nachricht an dch beschrieben wurden. Alles muss die Form SCHLAGWORT: {...} haben. Versuche es erneut. 
	  } 
	   } 
	   } 
	 & \\ 
 

    \theutterance \stepcounter{utterance}  

    & & & \multicolumn{4}{p{0.6\linewidth}}{\cellcolor[rgb]{1,0.85,0.72}{%
	\makecell[{{p{\linewidth}}}]{% 
	  \tt {\tiny [GM$\langle$B]}  
	 ABLEHNUNG: {'C98', 'C81', 'C10', 'B38', 'C00', 'C38', 'C62', 'A07', 'C08', 'A83', 'C76'}\\ \tt ARGUMENT: {'Die momentane Zusammensetzung könnte optimiert werden durch Austausch von Items mit niedrigerer Priorität gegen solche mit höherer Priorität, um den Gesamtaufwand effizienter zu nutzen.'} 
	  } 
	   } 
	   } 
	  \\ 
 

    \theutterance \stepcounter{utterance}  

    & & \multicolumn{4}{p{0.6\linewidth}}{\cellcolor[rgb]{0.9,0.9,0.9}{%
	\makecell[{{p{\linewidth}}}]{% 
	  \tt {\tiny [A$\langle$GM]}  
	 ABLEHNUNG: {'C98', 'C81', 'C10', 'B38', 'C00', 'C38', 'C62', 'A07', 'C08', 'A83', 'C76'}\\ \tt ARGUMENT: {'Die momentane Zusammensetzung könnte optimiert werden durch Austausch von Items mit niedrigerer Priorität gegen solche mit höherer Priorität, um den Gesamtaufwand effizienter zu nutzen.'} 
	  } 
	   } 
	   } 
	 & \\ 
 

    \theutterance \stepcounter{utterance}  

    & \multicolumn{4}{p{0.6\linewidth}}{\cellcolor[rgb]{0.8,1,0.9}{%
	\makecell[{{p{\linewidth}}}]{% 
	  \tt {\tiny [A$\rangle$GM]}  
	 VORSCHLAG: {'A03', 'C56', 'A07', 'C00', 'B38', 'C38', 'C81', 'C98'}\\ \tt ARGUMENT: {"Höchste Prioritätsgüter für mich sind enthalten. 'A83' und 'C76' wurden entfernt, da sie meine Wirtschaftlichkeit nicht erhöhen. 'A03' und 'C56' maximieren den Gesamtwert innerhalb des Aufwandlimits."} 
	  } 
	   } 
	   } 
	 & & \\ 
 

    \theutterance \stepcounter{utterance}  

    & & \multicolumn{4}{p{0.6\linewidth}}{\cellcolor[rgb]{0.9,0.9,0.9}{%
	\makecell[{{p{\linewidth}}}]{% 
	  \tt {\tiny [GM$\rangle$B]}  
	 VORSCHLAG: {'A03', 'C56', 'A07', 'C00', 'B38', 'C38', 'C81', 'C98'}\\ \tt ARGUMENT: {"Höchste Prioritätsgüter für mich sind enthalten. 'A83' und 'C76' wurden entfernt, da sie meine Wirtschaftlichkeit nicht erhöhen. 'A03' und 'C56' maximieren den Gesamtwert innerhalb des Aufwandlimits."} 
	  } 
	   } 
	   } 
	 & \\ 
 

    \theutterance \stepcounter{utterance}  

    & & & \multicolumn{4}{p{0.6\linewidth}}{\cellcolor[rgb]{1,0.85,0.72}{%
	\makecell[{{p{\linewidth}}}]{% 
	  \tt {\tiny [GM$\langle$B]}  
	 VORSCHLAG: {'A83', 'C08', 'C76', 'C98', 'C03', 'C10', 'C81', 'C62', 'C38', 'B38'}\\ \tt ARGUMENT: {"Diese Auswahl umfasst die höchsten Prioritäten für mich und sichert maximalen Nutzen innerhalb des erlaubten Aufwands."} 
	  } 
	   } 
	   } 
	  \\ 
 

    \theutterance \stepcounter{utterance}  

    & & \multicolumn{4}{p{0.6\linewidth}}{\cellcolor[rgb]{0.9,0.9,0.9}{%
	\makecell[{{p{\linewidth}}}]{% 
	  \tt {\tiny [A$\langle$GM]}  
	 VORSCHLAG: {'A83', 'C08', 'C76', 'C98', 'C03', 'C10', 'C81', 'C62', 'C38', 'B38'}\\ \tt ARGUMENT: {"Diese Auswahl umfasst die höchsten Prioritäten für mich und sichert maximalen Nutzen innerhalb des erlaubten Aufwands."} 
	  } 
	   } 
	   } 
	 & \\ 
 

    \theutterance \stepcounter{utterance}  

    & \multicolumn{4}{p{0.6\linewidth}}{\cellcolor[rgb]{0.8,1,0.9}{%
	\makecell[{{p{\linewidth}}}]{% 
	  \tt {\tiny [A$\rangle$GM]}  
	 VORSCHLAG: {'C98', 'C81', 'C10', 'B38', 'C00', 'C38', 'C62', 'A03', 'A07', 'C08'}\\ \tt ARGUMENT: {"Die Kombination optimiert die Nutzung des verfügbaren Aufwands mit einer größeren Anzahl hochwertiger Gegenstände für meine Prioritäten, während ineffiziente Elemente entfernt wurden."} 
	  } 
	   } 
	   } 
	 & & \\ 
 

    \theutterance \stepcounter{utterance}  

    & & \multicolumn{4}{p{0.6\linewidth}}{\cellcolor[rgb]{0.9,0.9,0.9}{%
	\makecell[{{p{\linewidth}}}]{% 
	  \tt {\tiny [GM$\rangle$B]}  
	 VORSCHLAG: {'C98', 'C81', 'C10', 'B38', 'C00', 'C38', 'C62', 'A03', 'A07', 'C08'}\\ \tt ARGUMENT: {"Die Kombination optimiert die Nutzung des verfügbaren Aufwands mit einer größeren Anzahl hochwertiger Gegenstände für meine Prioritäten, während ineffiziente Elemente entfernt wurden."} 
	  } 
	   } 
	   } 
	 & \\ 
 

    \theutterance \stepcounter{utterance}  

    & & & \multicolumn{4}{p{0.6\linewidth}}{\cellcolor[rgb]{1,0.85,0.72}{%
	\makecell[{{p{\linewidth}}}]{% 
	  \tt {\tiny [GM$\langle$B]}  
	 ABLEHNUNG: {'C98', 'C81', 'C10', 'B38', 'C00', 'C38', 'C62', 'A03', 'A07', 'C08'}\\ \tt ARGUMENT: {"'A03' hat für mich eine niedrige Wichtigkeit von 121, während 'A83' mit 868 fehlt. Ein Austausch könnte den Gesamtwert erheblich erhöhen."} 
	  } 
	   } 
	   } 
	  \\ 
 

    \theutterance \stepcounter{utterance}  

    & & \multicolumn{4}{p{0.6\linewidth}}{\cellcolor[rgb]{0.9,0.9,0.9}{%
	\makecell[{{p{\linewidth}}}]{% 
	  \tt {\tiny [A$\langle$GM]}  
	 ABLEHNUNG: {'C98', 'C81', 'C10', 'B38', 'C00', 'C38', 'C62', 'A03', 'A07', 'C08'}\\ \tt ARGUMENT: {"'A03' hat für mich eine niedrige Wichtigkeit von 121, während 'A83' mit 868 fehlt. Ein Austausch könnte den Gesamtwert erheblich erhöhen."} 
	  } 
	   } 
	   } 
	 & \\ 
 

    \theutterance \stepcounter{utterance}  

    & \multicolumn{4}{p{0.6\linewidth}}{\cellcolor[rgb]{0.8,1,0.9}{%
	\makecell[{{p{\linewidth}}}]{% 
	  \tt {\tiny [A$\rangle$GM]}  
	 VORSCHLAG: {'A83', 'C56', 'C00', 'B38', 'C38', 'C62', 'C81', 'C98', 'A07'}\\ \tt ARGUMENT: {"Diese Kombination enthält 'A83' wie gewünscht und optimiert den Gesamtwert durch die Aufnahme hochpriorisierter Items wie 'C56' und 'C00'."} 
	  } 
	   } 
	   } 
	 & & \\ 
 

    \theutterance \stepcounter{utterance}  

    & & \multicolumn{4}{p{0.6\linewidth}}{\cellcolor[rgb]{0.9,0.9,0.9}{%
	\makecell[{{p{\linewidth}}}]{% 
	  \tt {\tiny [GM$\rangle$B]}  
	 VORSCHLAG: {'A83', 'C56', 'C00', 'B38', 'C38', 'C62', 'C81', 'C98', 'A07'}\\ \tt ARGUMENT: {"Diese Kombination enthält 'A83' wie gewünscht und optimiert den Gesamtwert durch die Aufnahme hochpriorisierter Items wie 'C56' und 'C00'."} 
	  } 
	   } 
	   } 
	 & \\ 
 

    \theutterance \stepcounter{utterance}  

    & & & \multicolumn{4}{p{0.6\linewidth}}{\cellcolor[rgb]{1,0.85,0.72}{%
	\makecell[{{p{\linewidth}}}]{% 
	  \tt {\tiny [GM$\langle$B]}  
	 ABLEHNUNG: {'A83', 'C56', 'C00', 'B38', 'C38', 'C62', 'C81', 'C98', 'A07'}\\ \tt ARGUMENT: {"'C56' hat für mich nur 65 Wert, während hochpriorisierte Items wie 'C76' (808) und 'C08' (822) fehlen. Dies reduziert den Gesamtnutzen erheblich."} 
	  } 
	   } 
	   } 
	  \\ 
 

    \theutterance \stepcounter{utterance}  

    & & \multicolumn{4}{p{0.6\linewidth}}{\cellcolor[rgb]{0.9,0.9,0.9}{%
	\makecell[{{p{\linewidth}}}]{% 
	  \tt {\tiny [A$\langle$GM]}  
	 ABLEHNUNG: {'A83', 'C56', 'C00', 'B38', 'C38', 'C62', 'C81', 'C98', 'A07'}\\ \tt ARGUMENT: {"'C56' hat für mich nur 65 Wert, während hochpriorisierte Items wie 'C76' (808) und 'C08' (822) fehlen. Dies reduziert den Gesamtnutzen erheblich."} 
	  } 
	   } 
	   } 
	 & \\ 
 

    \theutterance \stepcounter{utterance}  

    & \multicolumn{4}{p{0.6\linewidth}}{\cellcolor[rgb]{0.8,1,0.9}{%
	\makecell[{{p{\linewidth}}}]{% 
	  \tt {\tiny [A$\rangle$GM]}  
	 VORSCHLAG: {'C00', 'B38', 'C38', 'A07', 'C32', 'C76', 'C08', 'C81', 'C10', 'C98'}\\ \tt ARGUMENT: {"Die Auswahl enthält die von Ihnen gewünschten Items 'C76' und 'C08' und maximiert den Gesamtwert unter Beachtung der Aufwandsgrenze."} 
	  } 
	   } 
	   } 
	 & & \\ 
 

    \theutterance \stepcounter{utterance}  

    & & \multicolumn{4}{p{0.6\linewidth}}{\cellcolor[rgb]{0.9,0.9,0.9}{%
	\makecell[{{p{\linewidth}}}]{% 
	  \tt {\tiny [GM$\rangle$B]}  
	 VORSCHLAG: {'C00', 'B38', 'C38', 'A07', 'C32', 'C76', 'C08', 'C81', 'C10', 'C98'}\\ \tt ARGUMENT: {"Die Auswahl enthält die von Ihnen gewünschten Items 'C76' und 'C08' und maximiert den Gesamtwert unter Beachtung der Aufwandsgrenze."} 
	  } 
	   } 
	   } 
	 & \\ 
 

    \theutterance \stepcounter{utterance}  

    & & & \multicolumn{4}{p{0.6\linewidth}}{\cellcolor[rgb]{1,0.85,0.72}{%
	\makecell[{{p{\linewidth}}}]{% 
	  \tt {\tiny [GM$\langle$B]}  
	 ABLEHNUNG: {'C00', 'B38', 'C38', 'A07', 'C32', 'C76', 'C08', 'C81', 'C10', 'C98'}\\ \tt ARGUMENT: {"'A83' mit höchster Wertigkeit von 868 fehlt, während 'C32' mit lediglich 138 Werten enthalten ist. Austausch von C32 gegen A83 erhöht den Gesamtwert erheblich ohne den Aufwand zu überschreiten."} 
	  } 
	   } 
	   } 
	  \\ 
 

    \theutterance \stepcounter{utterance}  

    & & \multicolumn{4}{p{0.6\linewidth}}{\cellcolor[rgb]{0.9,0.9,0.9}{%
	\makecell[{{p{\linewidth}}}]{% 
	  \tt {\tiny [A$\langle$GM]}  
	 ABLEHNUNG: {'C00', 'B38', 'C38', 'A07', 'C32', 'C76', 'C08', 'C81', 'C10', 'C98'}\\ \tt ARGUMENT: {"'A83' mit höchster Wertigkeit von 868 fehlt, während 'C32' mit lediglich 138 Werten enthalten ist. Austausch von C32 gegen A83 erhöht den Gesamtwert erheblich ohne den Aufwand zu überschreiten."} 
	  } 
	   } 
	   } 
	 & \\ 
 

    \theutterance \stepcounter{utterance}  

    & \multicolumn{4}{p{0.6\linewidth}}{\cellcolor[rgb]{0.8,1,0.9}{%
	\makecell[{{p{\linewidth}}}]{% 
	  \tt {\tiny [A$\rangle$GM]}  
	 VORSCHLAG: {'C00', 'B38', 'C38', 'A07', 'A83', 'C76', 'C08', 'C81', 'C10', 'C98'}\\ \tt ARGUMENT: {"Ich habe 'C32' durch 'A83' ersetzt, wie von Ihnen gewünscht. Der Gesamtaufwand beträgt 2568 und liegt unter der maximalen Grenze von 2984. Damit wird die Priorität von 'A83' berücksichtigt und die verbleibenden Gegenstände bieten einen fairen Wert für beide Seiten."} 
	  } 
	   } 
	   } 
	 & & \\ 
 

    \theutterance \stepcounter{utterance}  

    & & \multicolumn{4}{p{0.6\linewidth}}{\cellcolor[rgb]{0.9,0.9,0.9}{%
	\makecell[{{p{\linewidth}}}]{% 
	  \tt {\tiny [GM$\rangle$B]}  
	 VORSCHLAG: {'C00', 'B38', 'C38', 'A07', 'A83', 'C76', 'C08', 'C81', 'C10', 'C98'}\\ \tt ARGUMENT: {"Ich habe 'C32' durch 'A83' ersetzt, wie von Ihnen gewünscht. Der Gesamtaufwand beträgt 2568 und liegt unter der maximalen Grenze von 2984. Damit wird die Priorität von 'A83' berücksichtigt und die verbleibenden Gegenstände bieten einen fairen Wert für beide Seiten."} 
	  } 
	   } 
	   } 
	 & \\ 
 

    \theutterance \stepcounter{utterance}  

    & & & \multicolumn{4}{p{0.6\linewidth}}{\cellcolor[rgb]{1,0.85,0.72}{%
	\makecell[{{p{\linewidth}}}]{% 
	  \tt {\tiny [GM$\langle$B]}  
	 ABLEHNUNG: {'C00', 'B38', 'C38', 'A07', 'A83', 'C76', 'C08', 'C81', 'C10', 'C98'}\\ \tt ARGUMENT: {"'A07' mit Wichtigkeit 262 sollte durch 'C62' (Wichtigkeit 508) ersetzt werden, um den Gesamtwert zu erhöhen und den Aufwand effizient zu nutzen."} 
	  } 
	   } 
	   } 
	  \\ 
 

    \theutterance \stepcounter{utterance}  

    & & \multicolumn{4}{p{0.6\linewidth}}{\cellcolor[rgb]{0.9,0.9,0.9}{%
	\makecell[{{p{\linewidth}}}]{% 
	  \tt {\tiny [A$\langle$GM]}  
	 ABLEHNUNG: {'C00', 'B38', 'C38', 'A07', 'A83', 'C76', 'C08', 'C81', 'C10', 'C98'}\\ \tt ARGUMENT: {"'A07' mit Wichtigkeit 262 sollte durch 'C62' (Wichtigkeit 508) ersetzt werden, um den Gesamtwert zu erhöhen und den Aufwand effizient zu nutzen."} 
	  } 
	   } 
	   } 
	 & \\ 
 

    \theutterance \stepcounter{utterance}  

    & \multicolumn{4}{p{0.6\linewidth}}{\cellcolor[rgb]{0.8,1,0.9}{%
	\makecell[{{p{\linewidth}}}]{% 
	  \tt {\tiny [A$\rangle$GM]}  
	 VORSCHLAG: {'C00', 'B38', 'C38', 'A07', 'A83', 'C76', 'C08', 'C81', 'C10', 'C62'}\\ \tt ARGUMENT: {"'C98' wurde durch 'C62' ersetzt, um die Prioritäten beider Seiten zu optimieren, ohne den Maximalaufwand zu überschreiten. Diese Anpassung steigert den Nutzen für uns beide."} 
	  } 
	   } 
	   } 
	 & & \\ 
 

    \theutterance \stepcounter{utterance}  

    & & \multicolumn{4}{p{0.6\linewidth}}{\cellcolor[rgb]{0.9,0.9,0.9}{%
	\makecell[{{p{\linewidth}}}]{% 
	  \tt {\tiny [GM$\rangle$B]}  
	 VORSCHLAG: {'C00', 'B38', 'C38', 'A07', 'A83', 'C76', 'C08', 'C81', 'C10', 'C62'}\\ \tt ARGUMENT: {"'C98' wurde durch 'C62' ersetzt, um die Prioritäten beider Seiten zu optimieren, ohne den Maximalaufwand zu überschreiten. Diese Anpassung steigert den Nutzen für uns beide."} 
	  } 
	   } 
	   } 
	 & \\ 
 

    \theutterance \stepcounter{utterance}  

    & & & \multicolumn{4}{p{0.6\linewidth}}{\cellcolor[rgb]{1,0.85,0.72}{%
	\makecell[{{p{\linewidth}}}]{% 
	  \tt {\tiny [GM$\langle$B]}  
	 ABLEHNUNG: {'C00', 'B38', 'C38', 'A07', 'A83', 'C76', 'C08', 'C81', 'C10', 'C62'}\\ \tt ARGUMENT: {"C62 sollte durch C98 ersetzt werden, da C98 eine höhere Wichtigkeit von 783 für mich hat, während C62 nur 508 beträgt. Der Gesamtaufwand bleibt unter der Grenze."} 
	  } 
	   } 
	   } 
	  \\ 
 

    \theutterance \stepcounter{utterance}  

    & & \multicolumn{4}{p{0.6\linewidth}}{\cellcolor[rgb]{0.9,0.9,0.9}{%
	\makecell[{{p{\linewidth}}}]{% 
	  \tt {\tiny [A$\langle$GM]}  
	 ABLEHNUNG: {'C00', 'B38', 'C38', 'A07', 'A83', 'C76', 'C08', 'C81', 'C10', 'C62'}\\ \tt ARGUMENT: {"C62 sollte durch C98 ersetzt werden, da C98 eine höhere Wichtigkeit von 783 für mich hat, während C62 nur 508 beträgt. Der Gesamtaufwand bleibt unter der Grenze."} 
	  } 
	   } 
	   } 
	 & \\ 
 

    \theutterance \stepcounter{utterance}  

    & \multicolumn{4}{p{0.6\linewidth}}{\cellcolor[rgb]{0.8,1,0.9}{%
	\makecell[{{p{\linewidth}}}]{% 
	  \tt {\tiny [A$\rangle$GM]}  
	 assistant\\ \tt VORSCHLAG: {'C00', 'B38', 'C38', 'A07', 'C76', 'C08', 'C81', 'C10', 'C98', 'C56'}\\ \tt ARGUMENT: {"Die aktuelle Verteilung enthält hochpriorisierte Gegenstände für mich wie 'C56' (868) und 'C00' (780), während 'A83' entfernt wurde, da es für mich lediglich 65 Wert hat. Der Gesamtaufwand von 2795 liegt unter der Grenze und maximiert meinen Nutzen."} 
	  } 
	   } 
	   } 
	 & & \\ 
 

    \theutterance \stepcounter{utterance}  

    & & \multicolumn{4}{p{0.6\linewidth}}{\cellcolor[rgb]{0.9,0.9,0.9}{%
	\makecell[{{p{\linewidth}}}]{% 
	  \tt {\tiny [A$\langle$GM]}  
	 Deine Antwort hat eine unmarkierte Sequenz enthalten oder Du hast STRATEGIC REASONING mehr als ein mal verwendet. Du darfst nur die strukturierten Formate verwenden, die in der ersten Nachricht an dch beschrieben wurden. Alles muss die Form SCHLAGWORT: {...} haben. Versuche es erneut. 
	  } 
	   } 
	   } 
	 & \\ 
 

    \theutterance \stepcounter{utterance}  

    & \multicolumn{4}{p{0.6\linewidth}}{\cellcolor[rgb]{0.8,1,0.9}{%
	\makecell[{{p{\linewidth}}}]{% 
	  \tt {\tiny [A$\rangle$GM]}  
	 VORSCHLAG: {'C00', 'B38', 'C38', 'A07', 'C76', 'C08', 'C81', 'C10', 'C98', 'C56'}\\ \tt ARGUMENT: {"Die Auswahl enthält hochpriorisierte Gegenstände für mich, während weniger wichtige Items entfernt wurden. Der Gesamtaufwand liegt unter der Grenze und verbessert den Nutzen für beide Seiten."} 
	  } 
	   } 
	   } 
	 & & \\ 
 

    \theutterance \stepcounter{utterance}  

    & & \multicolumn{4}{p{0.6\linewidth}}{\cellcolor[rgb]{0.9,0.9,0.9}{%
	\makecell[{{p{\linewidth}}}]{% 
	  \tt {\tiny [GM$\rangle$B]}  
	 VORSCHLAG: {'C00', 'B38', 'C38', 'A07', 'C76', 'C08', 'C81', 'C10', 'C98', 'C56'}\\ \tt ARGUMENT: {"Die Auswahl enthält hochpriorisierte Gegenstände für mich, während weniger wichtige Items entfernt wurden. Der Gesamtaufwand liegt unter der Grenze und verbessert den Nutzen für beide Seiten."} 
	  } 
	   } 
	   } 
	 & \\ 
 

    \theutterance \stepcounter{utterance}  

    & & & \multicolumn{4}{p{0.6\linewidth}}{\cellcolor[rgb]{1,0.85,0.72}{%
	\makecell[{{p{\linewidth}}}]{% 
	  \tt {\tiny [GM$\langle$B]}  
	 ABLEHNUNG: {'C00', 'B38', 'C38', 'A07', 'C76', 'C08', 'C81', 'C10', 'C98', 'C56'}\\ \tt ARGUMENT: {"'C56' hat für mich eine niedrige Wichtigkeit von 65, während 'A83' mit höchster Wichtigkeit (868) fehlt. Ein Austausch von C56 durch A83 würde den Gesamtwert erheblich steigern."} 
	  } 
	   } 
	   } 
	  \\ 
 

    \theutterance \stepcounter{utterance}  

    & & \multicolumn{4}{p{0.6\linewidth}}{\cellcolor[rgb]{0.9,0.9,0.9}{%
	\makecell[{{p{\linewidth}}}]{% 
	  \tt {\tiny [A$\langle$GM]}  
	 ABLEHNUNG: {'C00', 'B38', 'C38', 'A07', 'C76', 'C08', 'C81', 'C10', 'C98', 'C56'}\\ \tt ARGUMENT: {"'C56' hat für mich eine niedrige Wichtigkeit von 65, während 'A83' mit höchster Wichtigkeit (868) fehlt. Ein Austausch von C56 durch A83 würde den Gesamtwert erheblich steigern."} 
	  } 
	   } 
	   } 
	 & \\ 
 

    \theutterance \stepcounter{utterance}  

    & \multicolumn{4}{p{0.6\linewidth}}{\cellcolor[rgb]{0.8,1,0.9}{%
	\makecell[{{p{\linewidth}}}]{% 
	  \tt {\tiny [A$\rangle$GM]}  
	 VORSCHLAG: {'A03', 'A83', 'C00', 'B38', 'A07', 'C62', 'C81', 'C98', 'C10'}\\ \tt ARGUMENT: {"Diese Auswahl enthält 'A83' wie gewünscht und schließt 'C56' aus. 'A03' und 'C00' maximieren meinen Nutzen, während der Gesamtaufwand von 2892 unter der Grenze liegt und beide Parteien gut bedient."} 
	  } 
	   } 
	   } 
	 & & \\ 
 

    \theutterance \stepcounter{utterance}  

    & & \multicolumn{4}{p{0.6\linewidth}}{\cellcolor[rgb]{0.9,0.9,0.9}{%
	\makecell[{{p{\linewidth}}}]{% 
	  \tt {\tiny [GM$\rangle$B]}  
	 VORSCHLAG: {'A03', 'A83', 'C00', 'B38', 'A07', 'C62', 'C81', 'C98', 'C10'}\\ \tt ARGUMENT: {"Diese Auswahl enthält 'A83' wie gewünscht und schließt 'C56' aus. 'A03' und 'C00' maximieren meinen Nutzen, während der Gesamtaufwand von 2892 unter der Grenze liegt und beide Parteien gut bedient."} 
	  } 
	   } 
	   } 
	 & \\ 
 

    \theutterance \stepcounter{utterance}  

    & & & \multicolumn{4}{p{0.6\linewidth}}{\cellcolor[rgb]{1,0.85,0.72}{%
	\makecell[{{p{\linewidth}}}]{% 
	  \tt {\tiny [GM$\langle$B]}  
	 ABLEHNUNG: {'A03', 'A83', 'C00', 'B38', 'A07', 'C62', 'C81', 'C98', 'C10'}\\ \tt ARGUMENT: {"'A03' hat eine niedrige Wichtigkeit von 121 für mich, während fehlende hochwertige Items wie 'C08' (822) und 'C76' (808) den Gesamtwert erheblich steigern würden."} 
	  } 
	   } 
	   } 
	  \\ 
 

    \theutterance \stepcounter{utterance}  

    & & \multicolumn{4}{p{0.6\linewidth}}{\cellcolor[rgb]{0.9,0.9,0.9}{%
	\makecell[{{p{\linewidth}}}]{% 
	  \tt {\tiny [A$\langle$GM]}  
	 ABLEHNUNG: {'A03', 'A83', 'C00', 'B38', 'A07', 'C62', 'C81', 'C98', 'C10'}\\ \tt ARGUMENT: {"'A03' hat eine niedrige Wichtigkeit von 121 für mich, während fehlende hochwertige Items wie 'C08' (822) und 'C76' (808) den Gesamtwert erheblich steigern würden."} 
	  } 
	   } 
	   } 
	 & \\ 
 

    \theutterance \stepcounter{utterance}  

    & \multicolumn{4}{p{0.6\linewidth}}{\cellcolor[rgb]{0.8,1,0.9}{%
	\makecell[{{p{\linewidth}}}]{% 
	  \tt {\tiny [A$\rangle$GM]}  
	 VORSCHLAG: {'C56', 'C00', 'A07', 'B38', 'C38', 'C62', 'C81', 'C76', 'C08'}\\ \tt ARGUMENT: {"Diese Auswahl enthält 'C08' und 'C76' wie gewünscht sowie hochpriorisierte Items wie 'C56' und 'C00', während weniger relevante Elemente entfernt wurden. Der Gesamtaufwand bleibt unter der Grenze von 2984."} 
	  } 
	   } 
	   } 
	 & & \\ 
 

    \theutterance \stepcounter{utterance}  

    & & \multicolumn{4}{p{0.6\linewidth}}{\cellcolor[rgb]{0.9,0.9,0.9}{%
	\makecell[{{p{\linewidth}}}]{% 
	  \tt {\tiny [GM$\rangle$B]}  
	 VORSCHLAG: {'C56', 'C00', 'A07', 'B38', 'C38', 'C62', 'C81', 'C76', 'C08'}\\ \tt ARGUMENT: {"Diese Auswahl enthält 'C08' und 'C76' wie gewünscht sowie hochpriorisierte Items wie 'C56' und 'C00', während weniger relevante Elemente entfernt wurden. Der Gesamtaufwand bleibt unter der Grenze von 2984."} 
	  } 
	   } 
	   } 
	 & \\ 
 

    \theutterance \stepcounter{utterance}  

    & & & \multicolumn{4}{p{0.6\linewidth}}{\cellcolor[rgb]{1,0.85,0.72}{%
	\makecell[{{p{\linewidth}}}]{% 
	  \tt {\tiny [GM$\langle$B]}  
	 ABLEHNUNG: {'C56', 'C00', 'A07', 'B38', 'C38', 'C62', 'C81', 'C76', 'C08'}\\ \tt ARGUMENT: {"'C56' hat eine niedrige Wichtigkeit von 65, während 'C98' mit 783 fehlt. Ein Austausch würde den Gesamtwert deutlich erhöhen, ohne den Aufwand zu überschreiten."} 
	  } 
	   } 
	   } 
	  \\ 
 

    \theutterance \stepcounter{utterance}  

    & & \multicolumn{4}{p{0.6\linewidth}}{\cellcolor[rgb]{0.9,0.9,0.9}{%
	\makecell[{{p{\linewidth}}}]{% 
	  \tt {\tiny [A$\langle$GM]}  
	 ABLEHNUNG: {'C56', 'C00', 'A07', 'B38', 'C38', 'C62', 'C81', 'C76', 'C08'}\\ \tt ARGUMENT: {"'C56' hat eine niedrige Wichtigkeit von 65, während 'C98' mit 783 fehlt. Ein Austausch würde den Gesamtwert deutlich erhöhen, ohne den Aufwand zu überschreiten."} 
	  } 
	   } 
	   } 
	 & \\ 
 

    \theutterance \stepcounter{utterance}  

    & \multicolumn{4}{p{0.6\linewidth}}{\cellcolor[rgb]{0.8,1,0.9}{%
	\makecell[{{p{\linewidth}}}]{% 
	  \tt {\tiny [A$\rangle$GM]}  
	 VORSCHLAG: {'C56', 'C00', 'A07', 'B38', 'C38', 'C62', 'C81', 'C98', 'C10'}\\ \tt ARGUMENT: {"Durch den Austausch von 'C76' gegen 'C98' und Hinzufügung von 'C10' wird Ihre Priorität für 'C98' berücksichtigt, während 'C56' erhalten bleibt. Der Gesamtaufwand beträgt 2658 und wirkt effizient für beide Seiten."} 
	  } 
	   } 
	   } 
	 & & \\ 
 

    \theutterance \stepcounter{utterance}  

    & & \multicolumn{4}{p{0.6\linewidth}}{\cellcolor[rgb]{0.9,0.9,0.9}{%
	\makecell[{{p{\linewidth}}}]{% 
	  \tt {\tiny [GM$\rangle$B]}  
	 VORSCHLAG: {'C56', 'C00', 'A07', 'B38', 'C38', 'C62', 'C81', 'C98', 'C10'}\\ \tt ARGUMENT: {"Durch den Austausch von 'C76' gegen 'C98' und Hinzufügung von 'C10' wird Ihre Priorität für 'C98' berücksichtigt, während 'C56' erhalten bleibt. Der Gesamtaufwand beträgt 2658 und wirkt effizient für beide Seiten."} 
	  } 
	   } 
	   } 
	 & \\ 
 

    \theutterance \stepcounter{utterance}  

    & & & \multicolumn{4}{p{0.6\linewidth}}{\cellcolor[rgb]{1,0.85,0.72}{%
	\makecell[{{p{\linewidth}}}]{% 
	  \tt {\tiny [GM$\langle$B]}  
	 assistant\\ \tt ABLEHNUNG: {'C56', 'C00', 'A07', 'B38', 'C38', 'C62', 'C81', 'C98', 'C10'}\\ \tt ARGUMENT: {"'C56' hat für mich eine Wichtigkeit von 65, während 'A83' (868) und 'C08' (822) fehlen. Ein Austausch von C56 durch A83 und Hinzufügung von C08 erhöht den Gesamtwert erheblich."} 
	  } 
	   } 
	   } 
	  \\ 
 

    \theutterance \stepcounter{utterance}  

    & & \multicolumn{4}{p{0.6\linewidth}}{\cellcolor[rgb]{0.9,0.9,0.9}{%
	\makecell[{{p{\linewidth}}}]{% 
	  \tt {\tiny [GM$\rangle$B]}  
	 Deine Antwort hat eine unmarkierte Sequenz enthalten oder Du hast STRATEGIC REASONING mehr als ein mal verwendet. Du darfst nur die strukturierten Formate verwenden, die in der ersten Nachricht an dch beschrieben wurden. Alles muss die Form SCHLAGWORT: {...} haben. Versuche es erneut. 
	  } 
	   } 
	   } 
	 & \\ 
 

    \theutterance \stepcounter{utterance}  

    & & & \multicolumn{4}{p{0.6\linewidth}}{\cellcolor[rgb]{1,0.85,0.72}{%
	\makecell[{{p{\linewidth}}}]{% 
	  \tt {\tiny [GM$\langle$B]}  
	 ABLEHNUNG: {"C56", "C00", "A07", "B38", "C38", "C62", "C81", "C98", "C10"}\\ \tt ARGUMENT: {"C56 hat für mich eine Wichtigkeit von 65, während A83 (868) und C08 (822) fehlen. Ein Austausch von C56 durch A83 und Hinzufügung von C08 erhöht den Gesamtwert erheblich."} 
	  } 
	   } 
	   } 
	  \\ 
 

    \theutterance \stepcounter{utterance}  

    & & \multicolumn{4}{p{0.6\linewidth}}{\cellcolor[rgb]{0.9,0.9,0.9}{%
	\makecell[{{p{\linewidth}}}]{% 
	  \tt {\tiny [A$\langle$GM]}  
	 ABLEHNUNG: {"C56", "C00", "A07", "B38", "C38", "C62", "C81", "C98", "C10"}\\ \tt ARGUMENT: {"C56 hat für mich eine Wichtigkeit von 65, während A83 (868) und C08 (822) fehlen. Ein Austausch von C56 durch A83 und Hinzufügung von C08 erhöht den Gesamtwert erheblich."} 
	  } 
	   } 
	   } 
	 & \\ 
 

    \theutterance \stepcounter{utterance}  

    & \multicolumn{4}{p{0.6\linewidth}}{\cellcolor[rgb]{0.8,1,0.9}{%
	\makecell[{{p{\linewidth}}}]{% 
	  \tt {\tiny [A$\rangle$GM]}  
	 VORSCHLAG: {'C56', 'C00', 'A07', 'B38', 'C38', 'C62', 'C81', 'C98', 'C10', 'C08'}\\ \tt ARGUMENT: {"C08 wurde hinzugefügt wie gewünscht, sodass jetzt sowohl C08 als auch C56 enthalten sind. Der Gesamtaufwand beträgt 2793 und liegt unter dem Limit. Dies optimiert die Werte für beide Seiten."} 
	  } 
	   } 
	   } 
	 & & \\ 
 

    \theutterance \stepcounter{utterance}  

    & & \multicolumn{4}{p{0.6\linewidth}}{\cellcolor[rgb]{0.9,0.9,0.9}{%
	\makecell[{{p{\linewidth}}}]{% 
	  \tt {\tiny [GM$\rangle$B]}  
	 VORSCHLAG: {'C56', 'C00', 'A07', 'B38', 'C38', 'C62', 'C81', 'C98', 'C10', 'C08'}\\ \tt ARGUMENT: {"C08 wurde hinzugefügt wie gewünscht, sodass jetzt sowohl C08 als auch C56 enthalten sind. Der Gesamtaufwand beträgt 2793 und liegt unter dem Limit. Dies optimiert die Werte für beide Seiten."} 
	  } 
	   } 
	   } 
	 & \\ 
 

    \theutterance \stepcounter{utterance}  

    & & & \multicolumn{4}{p{0.6\linewidth}}{\cellcolor[rgb]{1,0.85,0.72}{%
	\makecell[{{p{\linewidth}}}]{% 
	  \tt {\tiny [GM$\langle$B]}  
	 assistant\\ \tt ABLEHNUNG: {'C56', 'C00', 'A07', 'B38', 'C38', 'C62', 'C81', 'C98', 'C10', 'C08'}\\ \tt ARGUMENT: {"C56 hat für mich eine Wichtigkeit von 65, während A83 (868) fehlt. Ein Austausch von C56 durch A83 erhöht den Gesamtwert erheblich ohne den Aufwand zu überschreiten."} 
	  } 
	   } 
	   } 
	  \\ 
 

    \theutterance \stepcounter{utterance}  

    & & & \multicolumn{2}{p{0.3\linewidth}}{\cellcolor[rgb]{0.95,0.95,0.95}{%
	\makecell[{{p{\linewidth}}}]{% 
	  \tt {\tiny [GM$|$GM]}  
	 abort 
	  } 
	   } 
	   } 
	 & & \\ 
 

    \theutterance \stepcounter{utterance}  

    & & & \multicolumn{2}{p{0.3\linewidth}}{\cellcolor[rgb]{0.95,0.95,0.95}{%
	\makecell[{{p{\linewidth}}}]{% 
	  \tt {\tiny [GM$|$GM]}  
	 end game 
	  } 
	   } 
	   } 
	 & & \\ 
 

\end{supertabular}
}

\end{document}
